\documentclass[11pt]{article}
\usepackage[review]{acl}
\usepackage{times}
\usepackage{latexsym}
\usepackage{booktabs}
\usepackage[T1]{fontenc}
\usepackage[utf8]{inputenc}

\title{Section 5 Draft: Results}
\author{Anonymous ARR Submission}

\begin{document}
\maketitle

\section{Results}
\label{sec:results}

This section reports three sets of results: the full 300-instance SWE-bench Lite run, the fixed-slice ablation-labeled variants, and the controlled semantic-contract comparison.  Throughout, we distinguish official resolved rate from patch generation rate.

\subsection{Overall Benchmark Performance}
\label{subsec:overall-results}

Table~\ref{tab:overall-results} summarizes the main benchmark outcome.  Across the 300 submitted SWE-bench Lite instances, the official harness summaries report 58 resolved instances, giving a resolved rate of 19.3\%.  If the denominator is restricted to the 279 completed non-empty-patch instances reported by the evaluation summaries, the completed-only resolved rate is 20.8\%.

\begin{table}[t]
\centering
\small
\begin{tabular}{lr}
\toprule
Metric & Value \\
\midrule
Submitted instances & 300 \\
Completed non-empty-patch instances & 279 \\
Official resolved instances & 58 \\
Resolved / submitted & 19.3\% \\
Resolved / completed & 20.8\% \\
Patch generation rate (run logs) & $\sim$94\% \\
Advanced Analysis tokens & $\sim$28.3M \\
Mean runtime / instance & $\sim$523s \\
\bottomrule
\end{tabular}
\caption{Overall 300-instance benchmark summary.  Resolved rate is the official harness outcome; patch generation rate is a separate non-empty-diff measure.}
\label{tab:overall-results}
\end{table}

The difference between the 19.3\% resolved rate and the approximately 94\% patch generation rate is important.  The latter indicates that the workflow usually produces a syntactically applicable diff; it does not imply that the patch satisfies the benchmark tests.  THEMIS therefore appears better at producing candidate patches than at reliably producing correct issue resolutions under the current configuration.  This gap is consistent with the difficulty of repository-level issue repair, where a patch must localize the correct behavior, modify the correct file or symbol, and satisfy hidden or explicit tests.

The batch-level results are shown in Table~\ref{tab:batch-results}.  The strongest batch is cases 261--300, with 13 resolved instances out of 40 submitted.  Earlier batches have lower resolved counts, including 3/40 for cases 101--140.  We do not interpret these differences as repository-specific causal effects because the batches mix repositories and issue types.

\begin{table*}[t]
\centering
\small
\begin{tabular}{lrrrrr}
\toprule
Batch & Submitted & Completed & Empty patch & Resolved & Resolved / submitted \\
\midrule
cases 1--100 & 100 & 89 & 11 & 19 & 19.0\% \\
cases 101--140 & 40 & 38 & 2 & 3 & 7.5\% \\
cases 141--180 & 40 & 36 & 4 & 7 & 17.5\% \\
cases 181--220 & 40 & 39 & 1 & 8 & 20.0\% \\
cases 221--260 & 40 & 39 & 1 & 8 & 20.0\% \\
cases 261--300 & 40 & 38 & 2 & 13 & 32.5\% \\
\midrule
Total & 300 & 279 & 21 & 58 & 19.3\% \\
\bottomrule
\end{tabular}
\caption{Official evaluation summaries by batch.  The completed column corresponds to non-empty-patch instances in the evaluation summaries.}
\label{tab:batch-results}
\end{table*}

The full run covers the SWE-bench Lite Python repository set.  The largest approximate groups are Django ($\sim$70 instances), SymPy ($\sim$45), scikit-learn ($\sim$35), Matplotlib ($\sim$30), Sphinx ($\sim$20), Astropy ($\sim$18), Xarray ($\sim$16), Flask ($\sim$12), Seaborn ($\sim$12), Requests ($\sim$10), Pytest ($\sim$10), and other repositories ($\sim$22).  The run consumed approximately 28.3M tracked Advanced Analysis tokens, or about 94k tokens per instance, with an average wall-clock runtime of about 523 seconds per instance.

We exclude chat-model token counts from quantitative cost claims.  The logged \texttt{chat\_usage} fields are zero under the current callback path, but this is an accounting limitation rather than evidence that Developer or Judge calls were absent.

\subsection{Ablation-Labeled Variants}
\label{subsec:ablation-results}

The experiment log records six ablation-labeled variants on a fixed five-instance slice: \texttt{ablation1}, \texttt{ablation2}, \texttt{ablation13}, \texttt{ablation15}, \texttt{ablation15\_rankfix}, and \texttt{ablation456}.  The shared slice contains \texttt{django\_\_django-16139}, \texttt{sympy\_\_sympy-13773}, \texttt{django\_\_django-12497}, \texttt{django\_\_django-15320}, and \texttt{django\_\_django-12284}.  The documented output files provide patches for these variants, but the current planning materials do not define every numbered variant with enough precision to claim a fully isolated component effect.

We therefore treat this experiment as a design-probing slice rather than as a complete ablation proof.  Table~\ref{tab:ablation-variants} reports the verifiable setup information.  The clearest baseline is \texttt{ablation1}, which is documented as a graph-only conflict-detection baseline.  Other numbered variants should be interpreted cautiously until their exact component changes are defined in the manuscript or appendix.

\begin{table*}[t]
\centering
\small
\begin{tabular}{llll}
\toprule
Variant & Strategy / preset evidence & Instances & Interpretation boundary \\
\midrule
\texttt{ablation1} & \texttt{graph\_only}; graph-only conflict baseline & 5 & clearest structural baseline \\
\texttt{ablation2} & \texttt{graph\_only}; label documented & 5 & variant meaning requires definition \\
\texttt{ablation13} & label documented; details unspecified & 5 & do not over-attribute component effect \\
\texttt{ablation15} & label documented; details unspecified & 5 & do not over-attribute component effect \\
\texttt{ablation15\_rankfix} & \texttt{graph\_only}; ranking-fix label & 5 & ranking-related probe \\
\texttt{ablation456} & \texttt{graph\_only}; combined label & 5 & combined-variant probe \\
\bottomrule
\end{tabular}
\caption{Ablation-labeled variants on the fixed five-instance slice.  The table reports documented setup evidence and intentionally avoids unsupported causal claims for under-specified variant labels.}
\label{tab:ablation-variants}
\end{table*}

This conservative interpretation is useful for the paper.  It preserves the empirical record of the ablation-labeled runs while preventing the results section from implying that every architecture component has been independently isolated.  Detailed patch-level comparisons can be reported in an appendix once each variant is formally defined.

\subsection{Semantic Contract Analysis}
\label{subsec:semantic-contract-results}

The most interpretable controlled comparison is the four-case semantic-contract slice.  In this slice, the context-only control resolves 0/4 cases, the semantic-contract prompt resolves 1/4, and the reranking variant also resolves 1/4.  Table~\ref{tab:semantic-contract} summarizes the result.

\begin{table}[t]
\centering
\small
\begin{tabular}{lrr}
\toprule
Condition & Resolved & Total \\
\midrule
Context-only control & 0 & 4 \\
Semantic contract prompt & 1 & 4 \\
Semantic contract + rerank & 1 & 4 \\
\bottomrule
\end{tabular}
\caption{Controlled four-case semantic-contract slice.  Semantic contracts rescue one case relative to the context-only control, while reranking adds no observed gain in this slice.}
\label{tab:semantic-contract}
\end{table}

This result supports a bounded claim: explicit semantic contracts can help in at least some cases where additional context alone is insufficient.  However, the evidence is small and should not be read as a general statistical claim about all SWE-bench instances.  The negative reranking result is also informative.  It suggests that adding a more complex harness-aligned reranking step does not automatically improve over a clear semantic repair contract, at least in the observed slice.

The experiment log also records larger semantic-contract batches, including four-case, five-case, and seven-case semantic families with contract, control, and rerank outputs where available.  We use the four-case slice as the headline controlled result because it has the clearest documented control/contract/rerank comparison and aligns with the bounded claim in Section~\ref{sec:introduction}.

\bibliographystyle{acl_natbib}
\bibliography{../references}

\end{document}
